% Source: https://artofproblemsolving.com/community/c3h1951084_200s_post_problems
% Stable identity: aops-topic-1951084-q04
Write the number as \(\overline{d_1d_2d_3d_4}\) with \(1\le d_1\le 9\) and \(0\le d_2,d_3,d_4\le 9\). Put \(e_i=9-d_i\). Then \(0\le e_1\le 8\), \(0\le e_2,e_3,e_4\le 9\), and
\[
e_1+e_2+e_3+e_4=36-24=12.
\]
Without upper bounds there are \(\dbinom{15}{3}=455\) solutions. Subtract those violating a bound: \(e_1\ge 9\) leaves a sum of \(3\), giving \(\dbinom63=20\) solutions; \(e_i\ge 10\) for \(i=2,3,4\) leaves a sum of \(2\), giving \(\dbinom53=10\) each. Two violations at once are impossible (they would need a sum of at least \(19>12\)). Hence the count is
\[
455-20-3\cdot 10=\boxed{405}.
\]
